\chapter{The Language of Mathematics and the Art of Proof}

In a classroom, a student writes on the board: ``$n^{2}+n+41$ is prime for every
integer $n \geq 0$''. He then checks them one at a time. For $n=0$ he obtains
$41$; for $n=1$ he obtains $43$; for $n=2$ he obtains $47$. He continues up to
$n=39$, obtains $1601$, and $1601$ is indeed prime. Forty values, forty
failures that never came. At $n=40$ the claim breaks:
$40^{2}+40+41 = 1681 = 41^{2}$. One extra line cancels the whole preceding page.

That event is not an arithmetic accident but a grammatical mistake. ``For every
$n \geq 0$'' speaks about infinitely many numbers, whereas a list touches only
finitely many; no rule of logic moves truth from the finite to the infinite. This
chapter therefore offers a way of speaking such that checking is not needed.

The need is not a luxury. In Chapter~29 we will conclude that a
classifier can be encircled by shifting one input by $\varepsilon$,
and that conclusion holds for all inputs. In Chapter~27 a good accuracy report on
thousands of test data proves nothing about tomorrow's data. This chapter prepares
the apparatus for both.

\begin{untukanak}
On the shore of a lake lived a flock of geese. A child came every morning for
almost a year and counted the geese he saw: white, all of them, for two hundred
and seventy mornings. So he wrote in his notebook: all geese on this lake are white.

His grandfather asked how many geese the child had never seen. The child admitted he
did not know, and his grandfather said that the notebook told about the geese already
seen, not about all the geese on the lake.

This story is not about geese but about the difference between ``I have checked
everything'' and ``it holds for everything''. The first is only as strong as our memory. The second does not depend on who does the checking. One black goose
suffices to bring down a notebook of any thickness, whereas a thousand white geese
are never enough to set it up.

The black goose of this chapter is called forty. At that number the formula that looks tidy
suddenly yields a number that can be split. In this chapter we will learn why one
example brings down a claim, why many examples do not establish it, and how the gap
is closed without inspecting the last goose.
\end{untukanak}

\section{Statements, connectives, and truth tables}
\tingkat{basic}

A sentence can be judged only after we know what it promises. ``Seven is a
prime number'' promises something that can be checked, and that promise is true. ``Close
the window'' promises nothing. ``$x+1=3$'' cannot yet be judged, because its value
depends on $x$. The three look the same to spelling, but only the first
counts as mathematical material.

\begin{definisi}[Statement]
\index{statement}
A \emph{statement} is a sentence that is either true or false, but not
both at once.
\end{definisi}

A single statement becomes useful once it is combined. Each connective has one meaning
that is stipulated, not inferred from habits of speech.

\kamus{logical connective}{penghubung logika}{like a button on a calculator: two
lamps in, one lamp out, and the relation between the lamps is fixed at the factory,
not guessed by its user.}

With the picture of that button in mind, we stipulate the meaning of the five connectives officially.

\begin{definisi}[Logical connectives]
Let $P$ and $Q$ be statements.

We use the following five connectives.
\begin{itemize}
  \item $\neg P$ (\emph{negation}, \emph{not}): true exactly when $P$ is false.
  \item $P \wedge Q$ (\emph{conjunction}, \emph{and}): true exactly when both are true.
  \item $P \vee Q$ (\emph{disjunction}, \emph{or}): true exactly when at least
        one is true; the word ``or'' here is inclusive.
  \item $P \Rightarrow Q$ (\emph{implication}, \emph{if--then}): false exactly when
        $P$ is true and $Q$ is false; true in the other three cases.
  \item $P \Leftrightarrow Q$ (\emph{biconditional}, \emph{iff}): true exactly when
        $P$ and $Q$ have the same value.
\end{itemize}
\end{definisi}

All the content of that definition can be compressed into a single table: the rows are all combinations
of values, the columns are the value of each connective.

{\centering
\begin{tabular}{cc|ccccc}
\hline
$P$ & $Q$ & $\neg P$ & $P \wedge Q$ & $P \vee Q$ & $P \Rightarrow Q$ & $P \Leftrightarrow Q$ \\
\hline
T & T & F & T & T & T & T \\
T & F & F & F & T & F & F \\
F & T & T & F & T & T & F \\
F & F & T & F & F & T & T \\
\hline
\end{tabular}\par}

\begin{figure}[htbp]
\centering
\begin{tikzpicture}[font=\small]
  \fill[layerjade!16] (0,0) rectangle (3,1.8);
  \fill[accent!16] (3,0) rectangle (6,1.8);
  \fill[layerjade!16] (0,-1.8) rectangle (3,0);
  \fill[layerjade!16] (3,-1.8) rectangle (6,0);
  \draw[layergrey!70,thick] (0,-1.8) grid[xstep=3,ystep=1.8] (6,1.8);
  \draw[accent,very thick] (3,0) rectangle (6,1.8);
  \node at (1.5,0.9) {T};
  \node at (4.5,0.9) {F};
  \node at (1.5,-0.9) {T};
  \node at (4.5,-0.9) {T};
  \node[above,font=\footnotesize] at (1.5,1.8) {$Q$ true};
  \node[above,font=\footnotesize] at (4.5,1.8) {$Q$ false};
  \node[left,font=\footnotesize] at (0,0.9) {$P$ true};
  \node[left,font=\footnotesize] at (0,-0.9) {$P$ false};
  \node[below,font=\footnotesize,text=accent] at (4.5,-1.8) {the only false cell};
\end{tikzpicture}
\caption{The truth grid of $P \Rightarrow Q$: rows hold the value of $P$, columns hold the value of $Q$.
Three cells hold T and one cell holds F, namely when $P$ is true while $Q$ is false.
Anyone who wants to bring the implication down need only point at the cell framed in thick lines.}
\label{fig:kisi-implikasi}
\end{figure}

The third and fourth rows are the ones that most often surprise. In the fourth row,
$P$ is false and $Q$ is false, and the implication is judged true. In the third row, $P$ is false
and $Q$ is true, and the implication is true as well. So the only way to bring down
$P \Rightarrow Q$ is to exhibit $P$ true while $Q$ is false; the grid in
Figure~\ref{fig:kisi-implikasi} contains that cell as the only false cell.

That stipulation is not a matter of style. Consider ``if $n$ is divisible by $4$, then $n$
is even'', which must hold for every integer $n$, including $n=3$. There the
antecedent is false; if the implication were judged false, the statement that is plainly true would
collapse because of one irrelevant number. An implication records a relation of values, not
cause and effect.

\kamus{implication}{implikasi}{like the promise ``if it rains, the road is wet'': the promise is
broken only when rain falls and the road stays dry; when there is no rain, the one who made the
promise has done nothing wrong.}

Two compound statements can have the same value in every case; such an equality
has a name of its own.

\begin{definisi}[Tautology and equivalence]
\index{tautology}
A compound statement is called a \emph{tautology} if it is true under every assignment
of values to its constituent statements. Two compound statements $A$ and $B$ are called
\emph{logically equivalent}, written $A \equiv B$, if $A \Leftrightarrow B$ is a tautology.
\end{definisi}

For two constituent statements there are $2^{2}=4$ assignments of values; for three constituents there are
$2^{3}=8$. An equivalence that involves only two or three statements can therefore be
proved by checking every row, and that check is complete. This is the only place in this chapter where
exhaustive checking is accepted as proof, because the set of cases is finite and
has been counted.

\begin{proposisi}
For every statement $P$ and $Q$ we have $P \Rightarrow Q \equiv \neg P \vee Q$.
\end{proposisi}
\par\smallskip\noindent\textit{Proof.} We compare the column $P \Rightarrow Q$ with the column $\neg P \vee Q$
in the four rows of the truth table. In the first row both are true. In the second row $\neg P$ is false and
$Q$ is false, so both are false. In the third and fourth rows $\neg P$ is true,
so both are true. The four rows give the same value, hence the two statements are
equivalent.
\hfill$\blacksquare${\parfillskip=0pt\par}

Two derived equivalences are used throughout this book. The first lets us reverse
the direction of the talk without changing the value.

\begin{akibat}[Contraposition]
\index{contraposition}
For every statement $P$ and $Q$ we have
$P \Rightarrow Q \equiv \neg Q \Rightarrow \neg P$.
\end{akibat}
\par\smallskip\noindent\textit{Proof.} By the preceding proposition, both sides equal $\neg P \vee Q$ and
$\neg(\neg Q) \vee \neg P$. Since $\neg(\neg Q) \equiv Q$ and a disjunction pays no
attention to order, the two forms are identical.
\hfill$\blacksquare${\parfillskip=0pt\par}

What is \emph{not} equivalent is the converse, namely $Q \Rightarrow P$. In the third row of
the table, $Q \Rightarrow P$ is true while $P \Rightarrow Q$ is false. Swapping antecedent
for consequent is the most frequent mistake, and it is often invisible because the sentence
still reads smoothly.

\begin{proposisi}
For every statement $P$ and $Q$ we have
$\neg(P \Rightarrow Q) \equiv P \wedge \neg Q$.
\end{proposisi}
\par\smallskip\noindent\textit{Proof.} By the definition of negation and the first proposition,
$\neg(P \Rightarrow Q) \equiv \neg(\neg P \vee Q)$. De Morgan's law, proved in
Chapter~2, gives $\neg(\neg P \vee Q) \equiv \neg(\neg P) \wedge \neg Q$, and
$\neg(\neg P) \equiv P$. Hence $\neg(P \Rightarrow Q) \equiv P \wedge \neg Q$.
A reader who has not yet accepted De Morgan's law can check the four rows of the table
directly.
\hfill$\blacksquare${\parfillskip=0pt\par}

\begin{tangan}
Prepare eight small cards: four bearing T, four bearing F. Lay out four rows,
each row with one card for $P$ and one for $Q$: T-T, T-F, F-T, F-F. To the right of each
row place a third card that you write yourself: T if you judge ``if $P$ then $Q$'' to be
true, F if false.

Work without looking at the table, then compare. Only one row gets the F card, namely the
row whose $P$ is T and whose $Q$ is F. Slide that card up and down; it remains
the only one. Now swap the P and Q cards in the T-F row so that it reads F-T, then evaluate
again: the card changes from F to T. That one-card move is the converse.
\end{tangan}
\section{Quantifiers and how to negate them}
\tingkat{basic}

Most mathematical sentences speak about many objects, and that multiplicity is
marked by quantifiers. There are two quantifiers, and both read order as part of the meaning.

\kamus{quantifier}{kuantor}{like two kinds of announcement at school: ``all pupils
must bring a book'' and ``some pupil brings a book''; the two sound similar
but demand very different proofs.}

We write down their meaning officially.

\begin{definisi}[Universal and existential quantifiers]
\index{quantifier}
Let $P(x)$ be a statement that depends on $x$ inside the set $A$.
\begin{itemize}
  \item $(\forall x \in A)\,P(x)$ means $P(x)$ is true for every $x \in A$.
  \item $(\exists x \in A)\,P(x)$ means there is at least one $x \in A$
        with $P(x)$ true.
\end{itemize}
The first statement is false if there is one $x \in A$ with $P(x)$ false; one such
object is called a \emph{counterexample}.
\end{definisi}

The existential quantifier demands at least one, not exactly one; adding a demand
that is not written in the definition is the most frequent source of error.

The order of two quantifiers of different kinds may not be swapped. Inside $\mathbb{R}$,
the statement
\[
  (\forall x)(\exists y)\; y > x
\]
is true, whereas
\[
  (\exists y)(\forall x)\; y > x
\]
is false. The proofs of both are short.

\begin{proposisi}
Inside $\mathbb{R}$ we have $(\forall x)(\exists y)\, y > x$, but we do not have
$(\exists y)(\forall x)\, y > x$.
\end{proposisi}
\par\smallskip\noindent\textit{Proof.} For the first part, take any $x \in \mathbb{R}$
and choose $y = x+1$. Then $y > x$, so $y$ is a witness for that $x$.
Since $x$ was arbitrary, every $x$ has a witness; the witness may differ for different
$x$.

For the second part, suppose there is $y_{0} \in \mathbb{R}$ with $y_{0} > x$ for
every $x \in \mathbb{R}$. Take $x = y_{0}+1$, which is also a real number. Then
$y_{0} > y_{0}+1$, that is $0 > 1$, which is impossible. So no such $y_{0}$
exists.
\hfill$\blacksquare${\parfillskip=0pt\par}

\index{negation}Negation works from the outside in: every quantifier changes kind, and the negation travels inward
until it sits on the innermost statement. For compound statements, De Morgan's laws
move the negation past conjunctions and disjunctions.

{\centering
\begin{tabular}{ll}
\hline
Statement & Its negation \\
\hline
$(\forall x \in A)\,P(x)$ & $(\exists x \in A)\,\neg P(x)$ \\
$(\exists x \in A)\,P(x)$ & $(\forall x \in A)\,\neg P(x)$ \\
$P \wedge Q$ & $\neg P \vee \neg Q$ \\
$P \vee Q$ & $\neg P \wedge \neg Q$ \\
$P \Rightarrow Q$ & $P \wedge \neg Q$ \\
\hline
\end{tabular}\par}

The last row is the one most often violated. Denying ``if it rains then the road
is wet'' is not ``if it rains then the road is dry'', but ``it rains and the road is not
wet''. The same mistake appears when people deny a definition: the verb inside it is
negated, the quantifier is left intact.

\begin{proposisi}
For every statement $P(x)$ and $Q(x)$ we have
\[
  \neg(\forall x)\bigl(P(x) \Rightarrow Q(x)\bigr)
  \;\equiv\; (\exists x)\bigl(P(x) \wedge \neg Q(x)\bigr).
\]
\end{proposisi}
\par\smallskip\noindent\textit{Proof.} A universal quantifier is negated into an existential
quantifier over the negation of the body:
\[
  \neg(\forall x)R(x) \;\equiv\; (\exists x)\neg R(x).
\]
The body $R(x)$ here is $P(x) \Rightarrow Q(x)$, whose negation is $P(x) \wedge
\neg Q(x)$ by the proposition of the previous part. So the negation has the form
$(\exists x)(P(x) \wedge \neg Q(x))$.
\hfill$\blacksquare${\parfillskip=0pt\par}

That form sets a clear task when one wants to reject a quantified statement: exhibit
one object that satisfies the antecedent and fails the consequent.

The next two negations are used repeatedly; the first is the material of Chapter~4.

\begin{contoh}[Negation of the definition of a limit]
The definition of a limit reads: the sequence $(a_{n})$ tends to $L$ if
\[
  (\forall \varepsilon > 0)(\exists N \in \mathbb{N})(\forall n \geq N)\;
  |a_{n} - L| < \varepsilon .
\]
Negating from the outside in gives
\[
  (\exists \varepsilon > 0)(\forall N \in \mathbb{N})(\exists n \geq N)\;
  |a_{n} - L| \geq \varepsilon .
\]
The reading of the second sentence: there is one tolerance $\varepsilon$ that no
$N$ can repair, because there is always a term after $N$ that deviates by at least
$\varepsilon$. To prove that $(a_{n})$ does not tend to $L$, we must find
that $\varepsilon$, then for every $N$ find $n \geq N$ that violates it; $n$ may
depend on $N$.
\end{contoh}

The second negation concerns the existence of primes. Euclid proved that for
every natural number $n$ there is a prime larger than $n$:
\[
  (\forall n \in \mathbb{N})(\exists p)\bigl(p \text{ prime} \wedge p > n\bigr).
\]
Its negation is
\[
  (\exists n \in \mathbb{N})(\forall p)\bigl(p \text{ prime} \Rightarrow p \leq n\bigr),
\]
that is: there is one bound $n$ such that all primes do not exceed $n$. Denying
Euclid's claim demands such a bound, not merely failing to find a large prime.

\begin{untukanak}
Picture a class with thirty pupils. The teacher writes on the board: ``all pupils
bring a pencil.'' To prove that sentence true, the teacher must check thirty desks; to
prove it false, one pupil without a pencil is enough.

Now the teacher writes another sentence: ``some pupil brings a pencil.'' This time
the situation is reversed. One pupil with a pencil is already enough to prove
that sentence. But to prove it false, the teacher must check thirty desks
and make sure there is not a single pencil.

The two sentences only swap a word at the front, but the work of checking them swaps
entirely. Mathematicians are careful with ``all'' and ``some'' because both
count how many things must be checked.

There is one more trap. If some pupil brings two pencils, is ``some pupil brings
a pencil'' still true? Of course it is true: that sentence demands at least one,
not exactly one. In this chapter we will learn why a demand that is not written may not
be added on one's own.
\end{untukanak}

\section{Reading a definition literally}
\tingkat{intermediate}

A mathematical definition is not a short explanation of a word; it is a contract that
holds in both directions. If a number is called odd when it has the form $2k+1$ for
some integer $k$, we are bound to two things: every number of the form $2k+1$
must be called odd, and every number called odd must have the form $2k+1$.

To read a definition means to read a single sentence of the form ``$x$ is called $B$ if and only
if $C(x)$'', even though the words ``only if'' are rarely written; the phrase ``is called'' already
contains both directions. The first direction grants permission to use the name $B$, the second grants permission
to use the condition $C(x)$, and reversing either one wrecks the whole argument.

\kamus{necessary condition}{syarat perlu}{like the condition for entering a building: without a membership
card we may not enter, but having a membership card does not yet mean we are already
seated inside.}

Compare two statements about an integer $n$: ``if $6$ divides $n$ then $3$
divides $n$'' is true, whereas ``if $3$ divides $n$ then $6$ divides $n$'' is false,
with counterexample $n=3$. So ``$3$ divides $n$'' is a necessary condition for ``$6$ divides
$n$'', but not a sufficient one. A good definition names both, since only
in that way can the name and the condition be exchanged.

\begin{bacakeras}
Read it aloud as a dialogue for two voices, alternating, with one breath at every full stop.

First voice: ``If it rains, the road is wet.''
Second voice: ``The road is wet. So it rains.''
First voice: ``Not necessarily. The road is wet when it is hosed down, too.''
Second voice: ``Still wet. Still raining.''
First voice: ``What I promised was one direction only. Rain first, then wet. The reverse direction was never promised.''

Repeat that dialogue once more, faster, until the ear catches the difference. The first
sentence puts the rain in front; the second moves the rain to the back. The sound turns
around, the promise does not.

The line to repeat three times out loud: ``if it rains then it is wet, not if
it is wet then it rains.''
\end{bacakeras}

Now consider the definition used most often in this book.

\begin{definisi}[Prime number]
\index{prime number}
An integer $p$ is called \emph{prime} if $p \geq 2$ and the only positive divisors
of $p$ are $1$ and $p$ itself.
\end{definisi}

In everyday words: a prime is larger than one and cannot be split
into a product of two smaller numbers; the two clauses are tested separately.

\begin{proposisi}
The number $1$ is not prime.
\end{proposisi}
\par\smallskip\noindent\textit{Proof.} The definition demands $p \geq 2$. For $p = 1$
that condition fails, so $1$ does not satisfy the definition, whatever we can
say about its divisors. Note that $1$ in fact passes the second clause:
the only positive divisors of $1$ are indeed $1$ and $1$ itself. So what sets it aside
is the first clause, not the second.
\hfill$\blacksquare${\parfillskip=0pt\par}

This kind of event explains why a definition must be read to the last clause.
A reader who skips the first clause will conclude that $1$ is prime, and
the uniqueness of the factorization $12 = 2^{2}\cdot 3$ in Chapter~28 will be lost, because $12$ may
also be written $1 \cdot 2^{2} \cdot 3$.

A definition also determines how to deny membership. To show that a
number is not prime, we must find a positive divisor that is neither $1$ nor
the number itself. Here is one example that often escapes a quick check.

\begin{proposisi}
The number $91$ is not prime. More generally, the statement ``every odd number
that is not divisible by $3$ and not divisible by $5$ is prime'' is false.
\end{proposisi}
\par\smallskip\noindent\textit{Proof.} In fact $91 = 7 \cdot 13$, and both factors
lie between $1$ and $91$, so $91$ has a positive divisor other than $1$ and
itself. The number $91$ is odd, is not divisible by $3$ because $9+1 = 10$ is not a multiple
of $3$, and is not divisible by $5$ because its last digit is neither $0$ nor $5$. So $91$
is a counterexample to that general statement. For comparison, $97$ is indeed prime; among the integers $2$ and $99$ there are
exactly $25$ primes, a count produced by the program in the lab section.
\hfill$\blacksquare${\parfillskip=0pt\par}

\begin{definisi}[Perfect square]
An integer $n$ is called a \emph{perfect square} if there is an integer $k$
with $n = k^{2}$.
\end{definisi}

By this definition, $0$ is a perfect square, because $0 = 0^{2}$, and so is $1$,
because $1 = 1^{2}$. Meanwhile $2$ is not a perfect square, and that statement needs
a proof, not an admission that we failed to find $k$.

\begin{proposisi}
The number $2$ is not a perfect square.
\end{proposisi}
\par\smallskip\noindent\textit{Proof.} Suppose $2 = k^{2}$ for some $k \in
\mathbb{Z}$. Since $k^{2} = 2 > 0$, we have $k \neq 0$. If $|k| = 1$ then
$k^{2} = 1 \neq 2$. If $|k| \geq 2$ then $k^{2} \geq 4 > 2$. Since $|k|$ is a
positive integer, only those two possibilities remain, and both fail. Hence there is no
such $k$.
\hfill$\blacksquare${\parfillskip=0pt\par}

\index{counterexample}A habit worth keeping: every time you meet a
definition, ask three questions. What are its clauses? Which objects satisfy
all of them? Which objects are often thought to satisfy them although they fail one clause?
The third question produces counterexamples, and a counterexample can be neither
argued with nor debated.

\section{Four patterns of proof}
\tingkat{intermediate}

A proof is a transport of truth: from statements we have already accepted to the
statement we want to accept, through steps that can each be checked.
The four patterns below cover most needs, and
Figure~\ref{fig:peta-tiga-bukti} maps the first three routes. The fifth common pattern,
proof by cases, is merely the first pattern applied to cases that exclude one another.

\begin{figure}[htbp]
\centering
\begin{tikzpicture}[font=\small,>=Stealth,
  jalur/.style={draw,rounded corners=2pt,align=center,inner sep=5pt,text width=4.6cm},
  ujung/.style={draw,rounded corners=2pt,align=center,inner sep=4pt,text width=2.2cm}]
  \node[jalur,fill=layerorange!12] (klaim) at (0,0) {claim: $P \Rightarrow Q$};
  \node[jalur,fill=layerteal!12] (lgs) at (6.1,2.4)
    {\textbf{direct}\\ unpack $P$, compute, out comes $Q$};
  \node[jalur,fill=layerindigo!10] (ktp) at (6.1,0)
    {\textbf{contraposition}\\ prove $\neg Q \Rightarrow \neg P$; the value is the same};
  \node[jalur,fill=layerplum!10] (ktr) at (6.1,-2.4)
    {\textbf{contradiction}\\ suppose $P$ true and $Q$ false, derive a clash};
  \node[ujung,fill=layergrey!10] (selesai) at (11.6,0)
    {no assignment\\brings it\\down};
  \draw[->,thick] (klaim.east) -- (lgs.west);
  \draw[->,thick] (klaim.east) -- (ktp.west);
  \draw[->,thick] (klaim.east) -- (ktr.west);
  \draw[->,thick] (lgs.east) -- (selesai.north west);
  \draw[->,thick] (ktp.east) -- (selesai.west);
  \draw[->,thick] (ktr.east) -- (selesai.south west);
\end{tikzpicture}
\caption{Road map of three forms of proof for an implication. All routes start from the same claim
and end at the fact that no assignment of values makes $P$ true while
$Q$ is false. A proof by contraposition is valid because $P \Rightarrow Q \equiv \neg Q \Rightarrow \neg P$.
The fourth form, a constructive proof with a witness, does not enter this map because it proves
a statement of the form $\exists$, not an implication.}
\label{fig:peta-tiga-bukti}
\end{figure}

\paragraph{Pattern one: direct proof.} We start from the hypothesis, unpack it
by definition, then compute until the conclusion appears. The next example is chosen
because its conclusion is stronger than the one people usually prove.

\begin{teorema}
For every odd integer $n$, the number $n^{2}-1$ is divisible by $8$.
\end{teorema}
\par\smallskip\noindent\textit{Proof.} Since $n$ is odd, $n = 2k+1$ for some
$k \in \mathbb{Z}$. Then
\[
  n^{2} - 1 = (2k+1)^{2} - 1 = 4k^{2} + 4k = 4k(k+1).
\]
Among two consecutive integers $k$ and $k+1$, exactly one is even,
so $k(k+1) = 2m$ for some $m \in \mathbb{Z}$. Substitution gives
$n^{2}-1 = 4 \cdot 2m = 8m$, which means $n^{2}-1$ is divisible by $8$.
\hfill$\blacksquare${\parfillskip=0pt\par}

An immediate consequence: the square of an odd number always leaves remainder $1$ when divided by $8$. For
$n = 1,3,5,7,9,11$ we obtain in turn $n^{2}-1 = 0, 8, 24, 48, 80, 120$, and
all six are divisible by $8$. That check is not a proof, but it makes sure that
the proof above contains no typo.

\begin{figure}[htbp]
\centering
\begin{tikzpicture}[font=\small,>=Stealth,
  tahap/.style={draw,rounded corners=2pt,align=left,inner sep=5pt,text width=6.6cm},
  nama/.style={font=\footnotesize\itshape,text=layergrey,align=right,text width=2.4cm}]
  \node[tahap,fill=layerteal!10] (asumsi) at (0,0)
    {$n$ is odd, so $n = 2k+1$ for some $k \in \mathbb{Z}$};
  \node[nama] at (-4.9,0) {assumption (given)};
  \node[tahap,fill=layerorange!10] (langkah) at (0,-2.1)
    {$n^{2}-1 = 4k^{2}+4k = 4k(k+1) = 8m$};
  \node[nama] at (-4.9,-2.1) {step (computed)};
  \node[tahap,fill=layerjade!10] (kesimpulan) at (0,-4.2)
    {$8$ divides $n^{2}-1$};
  \node[nama] at (-4.9,-4.2) {conclusion (required)};
  \draw[->,thick] (asumsi.south) -- (langkah.north);
  \draw[->,thick] (langkah.south) -- (kesimpulan.north);
\end{tikzpicture}
\caption{Anatomy of one direct proof of the statement ``if $n$ is odd then $8$ divides
$n^{2}-1$''. The top box is the only one that may be used without reason; the middle box holds
its reason, namely the definition of an odd number and algebra; the bottom box is what is required.
The arrows may only run downward.}
\label{fig:anatomi-bukti}
\end{figure}

Figure~\ref{fig:anatomi-bukti} names the three parts of that proof: assumption, step,
conclusion.

\paragraph{Pattern two: proof by contraposition.} We prove $\neg Q \Rightarrow \neg P$
instead of $P \Rightarrow Q$. By the corollary in the first section the two are equivalent,
so this is valid. The pattern is chosen when the negation of the conclusion is easier to hold than
its hypothesis.

\begin{proposisi}
For every integer $n$, if $5n+3$ is even then $n$ is odd.
\end{proposisi}
\par\smallskip\noindent\textit{Proof.} The contrapositive reads: if $n$ is even then
$5n+3$ is odd. Let $n = 2k$. Then $5n+3 = 10k+3 = 2(5k+1)+1$, which has the form two
times an integer plus one, hence is odd. Since the contrapositive is proved, the
original implication is proved.
\hfill$\blacksquare${\parfillskip=0pt\par}

A direct proof would demand dividing $5n+3 = 2m$ by $5$; not impossible,
but longer. The pattern of contraposition is chosen because it saves work.

Before the third pattern, we need one lemma about parity. This lemma will be used twice:
once here, once in the proof that the square root of two is irrational.

\begin{lemma}
\index{parity lemma}
For every integer $n$, if $n^{2}$ is even then $n$ is even.
\end{lemma}
\par\smallskip\noindent\textit{Proof.} We prove the contrapositive: if $n$ is odd
then $n^{2}$ is odd. By the first theorem of this section, for odd $n$ we have
$n^{2} = 8m+1$ for some $m$, so $n^{2}$ leaves remainder $1$ when divided by $2$, that is,
is odd.
\hfill$\blacksquare${\parfillskip=0pt\par}

\paragraph{Pattern three: proof by contradiction.} We suppose the conclusion false, then derive two statements that cannot be true
together. This pattern is used when the object being denied has no convenient form
to compute with, for instance when what is denied is the existence of something.

\begin{teorema}
\index{square root of two}
There is no rational number $r$ with $r^{2} = 2$.
\end{teorema}
\par\smallskip\noindent\textit{Proof.} Suppose there is $r \in \mathbb{Q}$ with
$r^{2} = 2$. Every rational number can be written as $p/q$ with $p,q$ integers,
$q \neq 0$, and $p,q$ having no common factor other than $1$; we
remove that common factor by dividing both, so a form without a common
factor can always be reached.

From $p^{2}/q^{2} = 2$ we obtain $p^{2} = 2q^{2}$, so $p^{2}$ is even. By
the lemma above, $p$ is even; write $p = 2m$ with $m \in \mathbb{Z}$. Substitution gives
$4m^{2} = 2q^{2}$, that is $q^{2} = 2m^{2}$, so $q^{2}$ is even and, once again
by the lemma, $q$ is even. Now $p$ and $q$ are both even, so both are divisible by
$2$, contradicting the choice of $p/q$ without a common factor. The supposition is false,
hence there is no rational $r$ with $r^{2} = 2$.
\hfill$\blacksquare${\parfillskip=0pt\par}

The contradiction is not between $r^{2}=2$ and ordinary arithmetic, but between ``$p,q$
both even'' and the opening promise ``$p,q$ without a common factor''. That is why the step
of removing common factors is not decoration; it is part of the hypothesis that is later violated.

\paragraph{Pattern four: constructive proof with a witness.} To prove a statement
of the form $\exists$, we point at one object and check that the object satisfies
the condition. The object pointed at is called a \emph{witness}.

\begin{teorema}
For every natural number $n$ there are $n$ consecutive integers that are all
not prime.
\end{teorema}
\par\smallskip\noindent\textit{Proof.} Take $N = (n+1)!$ and consider the $n$ numbers
\[
  N+2,\; N+3,\; \dots,\; N+(n+1).
\]
For every $k$ with $2 \leq k \leq n+1$, the number $k$ divides $(n+1)!$ because $k$
is one of its factors, so $k$ also divides $N+k$. Moreover
$N+k > k > 1$, so $N+k$ has a positive divisor $k$ that is neither $1$ nor
itself. By the definition, $N+k$ is not prime. We have pointed at the $n$ numbers
required, and there is no need to look for the smallest ones.
\hfill$\blacksquare${\parfillskip=0pt\par}

For $n = 5$ that proof points at $6! = 720$, then $722, 723, 724, 725, 726$. All five
indeed split: $722 = 2 \cdot 19^{2}$, $723 = 3 \cdot 241$, $724 = 2^{2} \cdot 181$,
$725 = 5^{2} \cdot 29$, and $726 = 2 \cdot 3 \cdot 11^{2}$. The number $721 = 7 \cdot 103$ splits as well, but it is not on our list.

A constructive proof does not promise the smallest or the fastest witness; what is judged is only
whether the object pointed at satisfies the condition. Denying a statement of the form $\exists$
demands showing that no object satisfies it, and that is where a contradiction
is needed.

\begin{resep}
Four standard steps to prove ``if $P$ then $Q$''.
\begin{enumerate}
  \item Write $P$ and $Q$ by their definitions, one line each.
  \item Open with ``suppose $P$ is true''. Only $P$ may be supposed; $Q$ may not
        yet be used.
  \item Derive $Q$ by substitution and algebra, one step per line, and name
        the definitions used.
  \item Close with a sentence that states $Q$ exactly as it appears in step 1.
\end{enumerate}
A warning about order: step 2 may not be swapped with step 3. Scratch work from the conclusion
backward is allowed, but the proof submitted must run from $P$ to $Q$; otherwise
what is proved is the converse.
\end{resep}
\section{Induction and strong induction}
\tingkat{intermediate}

Induction is a proof pattern for statements that come in a row, one for each natural
number. We set the first card upright and make sure that every card that falls knocks down
the next card; after both, no card may remain standing.

\kamus{induction}{induksi}{like a row of dominoes: it is enough to make sure the first card falls
and that every card knocks down the card in front of it, and then the whole row falls without needing
to be pushed one by one.}

We take the principle of induction as a basic property of $\mathbb{N} = \{1,2,3,\dots\}$, equivalent
to the statement that every non-empty subset of $\mathbb{N}$ has a
least member; nothing is proved circularly, we merely choose a starting point as in Chapter~3.

\begin{catatan}[Principle of induction]
Let $P(n)$ be a statement about natural numbers. If $P(1)$ is true and, for every
$k \in \mathbb{N}$, the truth of $P(k)$ forces the truth of $P(k+1)$, then $P(n)$ is true
for every $n \in \mathbb{N}$.
\end{catatan}

With that principle, we prove the formula that was already guessed in the opening of this section.

\begin{teorema}
For every natural number $n$ we have
\[
  1 + 3 + 5 + \cdots + (2n-1) = n^{2}.
\]
\end{teorema}
\par\smallskip\noindent\textit{Proof.} Base: for $n = 1$ the left side equals $1$ and
the right side equals $1^{2} = 1$. Step: suppose the formula is true for $n = k$, that is
$1+3+\cdots+(2k-1) = k^{2}$. The next term is $2(k+1)-1 = 2k+1$, so
\[
  1+3+\cdots+(2k-1)+(2k+1) = k^{2} + 2k + 1 = (k+1)^{2}.
\]
Hence the formula is true for $k+1$. By the principle of induction, the formula is true for every
$n \in \mathbb{N}$.
\hfill$\blacksquare${\parfillskip=0pt\par}

The second example has the form of an inequality. From $2^{k} \geq k+1$ and the condition $k \geq 0$
we obtain $2^{k+1} = 2 \cdot 2^{k} \geq 2(k+1) \geq k+2$, so with base $n = 0$
we have $2^{n} \geq n+1$ for every $n \geq 0$; the hypothesis and the side condition are used
at once.

Ordinary induction demands one step back. Sometimes the conclusion for $n$ follows only from
many earlier values at once, and there ordinary induction is not enough.

\begin{catatan}[Principle of strong induction]
If for every $n \in \mathbb{N}$ the truth of $P(1), P(2), \dots, P(n-1)$ forces
the truth of $P(n)$, then $P(n)$ is true for every $n \in \mathbb{N}$.
\end{catatan}

\kamus{strong induction}{induksi kuat}{like building a tower of blocks: to
place the tenth block we may use every block beneath it, not only
the ninth block.}

With that permission we can prove something ordinary induction cannot reach.

\begin{teorema}
Every integer $n \geq 2$ can be written as a product of prime
numbers.
\end{teorema}
\par\smallskip\noindent\textit{Proof.} We use strong induction on $n$. For $n = 2$:
the number $2$ is itself prime, so it is a product with one factor. Now take
$n \geq 3$ and suppose the statement true for all integers $m$ with
$2 \leq m < n$. If $n$ is prime, we are done. If $n$ is not prime, by the definition there is a
positive divisor $a$ with $a \neq 1$ and $a \neq n$; write $n = ab$ with
$b = n/a$. Since $a \neq 1$ and $a \neq n$, we have $2 \leq a \leq n-1$, and likewise
$2 \leq b \leq n-1$. The induction hypothesis applies to $a$ and to $b$, so both are
products of primes, and the product of both is $n$.
\hfill$\blacksquare${\parfillskip=0pt\par}

The strong induction hypothesis is used twice, on $a$ and on $b$. That is where ordinary induction
is awkward: both may be much smaller than $n-1$, so the statement for
$n-1$ alone does not cover them.

That theorem closes the promise of the previous section: a list of primes is never complete.

\begin{akibat}[Infinitely many primes]
\index{infinitely many primes}
For every finite list of primes $p_{1}, p_{2}, \dots, p_{k}$ there is
a prime that is not on that list.
\end{akibat}
\par\smallskip\noindent\textit{Proof.} Let $N = p_{1}p_{2}\cdots p_{k} + 1$. Since
$N \geq 2$, the previous theorem gives that $N$ is a product of primes;
take one prime factor $q$ of $N$. Suppose $q$ is on the list, say
$q = p_{i}$. Then $q$ divides the product $p_{1}\cdots p_{k}$, and since $q$ also
divides $N$, the number $q$ divides the difference $N - p_{1}\cdots p_{k} = 1$. No
prime divides $1$, so the supposition is false. Hence $q$ is not on
the list.
\hfill$\blacksquare${\parfillskip=0pt\par}

That proof does not prove that $N$ is prime, only that $N$ has a prime factor outside the
list: for $2,3,5,7$ we obtain $N = 211$, which happens to be prime, for $2,3,5,7,11,13$
we obtain $N = 30031 = 59 \cdot 509$. The proof is constructive; ``suppose there were only
finitely many primes'' is merely its rhetorical wrapping.

\begin{untukanak}
Picture a row of a thousand boxes, and we want to prove that every box contains
a marble. We do not need to open a thousand boxes. It is enough to open the first box, then prove
one thing: if a box contains a marble, the box beside it contains a marble too.
If both things are true, a thousand boxes are checked without being opened one by one.

This is induction, and it feels like magic. The key lies in the second promise: if the relation
between boxes leaks in even one place, the whole row collapses.

Now picture the first box turning out empty. The relation between boxes remains true, but
it does not help, because there is nothing first to fall. This is the mistake that
happens most often.

In this chapter we will learn why the two steps may not be swapped, and why
checking many boxes is no substitute for either of them.
\end{untukanak}

\section{When a proof fails}
\tingkat{advanced}

Back to the blackboard from the opening of the chapter. The first forty values are prime, and at
$n = 40$ the pattern breaks:
\[
  40^{2} + 40 + 41 = 1600 + 40 + 41 = 1681 = 41^{2}.
\]
That failure is no coincidence, and its shape can be explained without a computer. For
every integer $p \geq 2$, the substitution $n = p-1$ gives
\[
  (p-1)^{2} + (p-1) + p = p^{2} - 2p + 1 + p - 1 + p = p^{2},
\]
whereas the substitution $n = p$ gives $p^{2} + p + p = p(p+2)$. So every number
$p \geq 2$ produces two points of failure, namely $n = p-1$ and $n = p$. With $p = 41$
both produce $1681 = 41^{2}$ and $1763 = 41 \cdot 43$.

The same pattern appears in smaller polynomials. For $n^{2}+n+11$ the first point of failure
is $n = 10$, which gives $121 = 11^{2}$; for $n^{2}+n+17$ the first point of failure
is $n = 16$, which gives $289 = 17^{2}$; for $n^{2}+n+29$ the failure comes at $n = 2$,
because $35 = 5 \cdot 7$ is not prime. The larger the number in the constant position,
the longer the clean initial stretch; that is the whole secret of the reputation of $41$.

That failure yields three lessons that are often confused.

The first lesson is asymmetric. One counterexample brings down a statement
of the form ``for all'', whereas forty successful examples prove nothing
about the forty-first value. A computer check is therefore a tool for finding conjectures,
not a tool for proving.

The second lesson concerns the location of the defect. The first defect that most often escapes notice appears
in an induction proof. Consider the following false claim: for every $n \geq 1$
we have $1+3+\cdots+(2n-1) = n^{2}+1$. Its induction step, when inspected, turns out to be
valid:
\[
  \bigl(n^{2}+1\bigr) + (2n+1) = n^{2} + 2n + 2 = (n+1)^{2} + 1 .
\]
So from the truth of the formula at $n$ the truth at $n+1$ does indeed flow. What fails
is the base: for $n = 1$ the left side equals $1$ while the right side equals $2$.
The induction step moves truth, it does not create it.

The third lesson concerns the limits of checking. The usual divisibility test: $m > 1$ is prime if
it is not divisible by any prime smaller than $\sqrt{m}$. For $m = 1681$ we get $\sqrt{1681} = 41$, and the primes
smaller than $41$ are $2, 3, 5, 7, 11, 13, 17, 19, 23, 29, 31, 37$; not
one of them divides $1681$. By that rule $1681$ passes as prime. The error lies
in the word ``smaller'': the correct rule uses ``smaller than or equal to'',
because $1681 = 41 \cdot 41$ and the divisor is exactly the size of the root. One loose inequality sign
produces one false prime, precisely at $n = 40$.

Three other defects recur in longer writing. The defect of circular reasoning:
``since $n^{2}+n+41$ is prime for all $n$, it is also prime for $n = 40$'' uses
what must be proved as its reason. The defect of dividing by zero: from $a = b$ people
conclude $a+b = b$ after dividing by $a-b$, although $a-b = 0$. The defect of swapping
direction: from $P \Rightarrow Q$ and $Q$ people conclude $P$; in the truth table,
the third row is where that error hides.

All four defects share one root: a step is used outside its condition.
A useful habit is not to memorize patterns but to check the condition of every step:
in a contradiction the supposition must be the negation of the conclusion, in induction there is a base and a
step, in a division the divisor is not zero.

\begin{tantangan}
Claim: for every integer $k \geq 0$, the number $2^{2^{k}}+1$ is prime. For
$k = 0,1,2,3,4$ the values are $3$, $5$, $17$, $257$, $65537$, and all five are indeed prime.
Find the smallest $k$ that brings the claim down, then exhibit a divisor; the claim is
false, and the error can be shown without a computer.

Hint 1: five values are not a proof. Check $k = 5$; the number is $4294967297$, and its smallest
divisor is far smaller than its square root.

Hint 2: that divisor has the form $5 \cdot 2^{7}+1 = 641$, and
$4294967297 = 641 \cdot 6700417$.

If you get stuck, go back to the section ``When a proof fails''.
\end{tantangan}

\section{Worked examples}
\tingkat{intermediate}

The three computations below support this chapter and can be repeated on paper. The first is
the sum of the first odd numbers, term by term from $n = 1$.

{\centering
\begin{tabular}{cccc}
\hline
$n$ & terms & sum & $n^{2}$ \\
\hline
$1$ & $1$ & $1$ & $1$ \\
$2$ & $1+3$ & $4$ & $4$ \\
$3$ & $1+3+5$ & $9$ & $9$ \\
$4$ & $1+3+5+7$ & $16$ & $16$ \\
$5$ & $1+3+5+7+9$ & $25$ & $25$ \\
$6$ & $1+3+5+7+9+11$ & $36$ & $36$ \\
\hline
\end{tabular}\par}

The continuation: $36+13 = 49 = 7^{2}$, then $49+15 = 64$, $64+17 = 81$, $81+19 = 100$.
The difference of two consecutive squares is indeed $(n+1)^{2}-n^{2} = 2n+1$, and that is where induction
takes its stand.

The second is a primality test on $n^{2}+n+41$, only at the interesting values.

{\centering
\begin{tabular}{ccc}
\hline
$n$ & $n^{2}+n+41$ & remark \\
\hline
$0$ & $41$ & prime \\
$1$ & $43$ & prime \\
$39$ & $1601$ & prime \\
$40$ & $1681 = 41^{2}$ & not prime \\
$41$ & $1763 = 41 \cdot 43$ & not prime \\
\hline
\end{tabular}\par}

An exhaustive check for $n = 0$ up to $n = 40$ gives $40$ prime values out of $41$;
the only one that fails is $n = 40$, and $\sqrt{1681} = 41$ exactly, so the decisive
divisor sits right at the boundary of the check.

The third is the contraposition check on the truth table, written out in full.

{\centering
\begin{tabular}{cc|cc|c}
\hline
$P$ & $Q$ & $P \Rightarrow Q$ & $\neg Q \Rightarrow \neg P$ & equal \\
\hline
T & T & T & T & yes \\
T & F & F & F & yes \\
F & T & T & T & yes \\
F & F & T & T & yes \\
\hline
\end{tabular}\par}

The four rows give equality: the contrapositive pattern is not a rhetorical trick but a replacement
of form that does not change the value. The column $Q \Rightarrow P$ would differ in the third row,
and that difference is what makes the converse unusable.

\kertas{Compute $1+3+\cdots+(2n-1)$ for $n = 11$ and $n = 12$ by starting from
$100 = 10^{2}$, that is by adding the next two odd numbers one at a time. Then
compute $n^{2}+n+17$ for $n = 16$ and check by division whether $289$ equals
$17^{2}$. Finally, find the first point of failure of $n^{2}+n+11$ by testing
$n = 0, 1, 2, \dots$ and explain why that point can be predicted from the shape of
its polynomial.}

\section{Computational lab}
\tingkat{advanced}

The following script uses only the Python standard library and can be rerun without
any installation; every output line is genuine output.

\begin{lab}
\begin{verbatim}
from itertools import product

def prime(n):
    if n < 2:
        return False
    d = 2
    while d * d <= n:
        if n % d == 0:
            return False
        d += 1
    return True

pairs = product([False, True], repeat=2)
print("contraposition:", all((not P or Q) == (Q or not P)
                            for P, Q in pairs))

D = range(1, 11)
print("y>x:", sum(y > x for x in D for y in D), "of 100")
print("x without a witness:", [x for x in D if all(y <= x for y in D)])

values = [(n, n*n + n + 41) for n in range(41)]
print("prime:", sum(prime(v) for _, v in values), "of 41")
print("failed:", [(n, v, int(v**0.5)) for n, v in values if not prime(v)])
print("41*41, 41*43:", 41*41, 41*43)

residues = {(n*n - 1) % 8 for n in range(1, 40, 2)}
print("residues mod 8:", residues)

step = all(n*n + 1 + 2*n + 1 == (n+1)**2 + 1 for n in range(1, 500))
print("false step:", step)
\end{verbatim}
\end{lab}

The output obtained:

\begin{verbatim}
contraposition: True
y>x: 45 of 100
x without a witness: [10]
prime: 40 of 41
failed: [(40, 1681, 41)]
41*41, 41*43: 1681 1763
residues mod 8: {0}
false step: True
\end{verbatim}

The first line checks all four rows of the table at once. The second and third lines compute
what is awkward to do on paper: out of a hundred pairs $(x,y)$ in
$\{1,\dots,10\}$, $45$ satisfy $y > x$, and the only $x$ without a witness is $10$;
that is the failure of $(\forall x)(\exists y)$ on a finite set.

The fourth and fifth lines give $40$ primes out of $41$ values, failing only at
$(40,1681,41)$; that last $41$ is $\sqrt{1681}$ truncated. The sixth line
confirms $1681 = 41^{2}$ and $1763 = 41 \cdot 43$. The seventh states that every odd
$n$ below $40$ gives residue $0$ modulo $8$, in agreement with the theorem of the first pattern. The last
line is a warning: the induction step of that false claim is valid for the first five hundred values,
so a numerical check does not find its defect.

\begin{untukdosen}
A teaching order that has worked: start from $n = 40$ and $1681$, not from the
definitions. After the students are annoyed that forty checks were in vain, then
enter the truth table, with emphasis on the third and fourth rows. Continue to
the negation of quantifiers through the definition of the limit of a sequence, then to literal definitions with the example
$91 = 7 \cdot 13$; close with induction, a failing base, and the lab.

Two misconceptions are most frequent. First, an implication is read as cause and effect, so a row
with a false antecedent is judged false and a vacuous case proof is rejected. Second, induction
is thought finished once the step is proved, and strong induction is thought to be ordinary induction
repeated; as a result the base is not checked and the hypothesis is used without naming
its range.

Blackboard activity: split the class into two groups. The first group writes a quantified claim,
the second group writes its negation, then each names one object that
refutes the other side's claim. The game stops when a negation can no longer be refuted.

Students are ready to go on if they can negate the definition of a limit without opening their notes
and can explain why the base in induction may not be skipped.
\end{untukdosen}

\section{Summary}
\tingkat{basic}

\begin{ringkasan}
\begin{itemize}
  \item A statement is true or false; $P \Rightarrow Q$ is false only when $P$ is
        true and $Q$ is false.
  \item $P \Rightarrow Q \equiv \neg Q \Rightarrow \neg P$, but the converse is not
        equivalent; $\neg(P \Rightarrow Q) \equiv P \wedge \neg Q$.
  \item Negation works from the outside in: every quantifier changes kind, so
        denying ``for all'' means pointing at one counterexample.
  \item A definition holds in both directions: the name may be used only if every clause
        is satisfied, and every clause is used to check it.
  \item Four proof patterns: direct, contraposition, contradiction, constructive. Choose
        the shortest, not the most magnificent.
  \item Induction demands a base and a step; strong induction demands all earlier
        values. A step without a base proves nothing.
\end{itemize}
\end{ringkasan}

\section{Exercises}
\tingkat{intermediate}

\begin{latihan}
\begin{enumerate}
  \item Write the negation of the following statement, then turn the result into a
        plain sentence that states what must be shown in order to
        refute it: ``for every $\varepsilon > 0$ there is a natural number $N$
        such that for every $n \geq N$ we have $|a_{n} - L| < \varepsilon$''.

  \item Write the negation of ``every prime larger than $2$ is
        odd'', then name one object that refutes the original statement.
        Do the same for ``there is a natural number $n$ such that
        $n^{2}+n+41$ is not prime''.

  \item Prove that for every integer $n$, the number $n^{2}+n$ is always
        even. Do it in two ways: by cases on parity, and by contraposition.

  \item Prove by induction that for every natural number $n$ we have
        \[
          1^{2} + 2^{2} + \cdots + n^{2} = \frac{n(n+1)(2n+1)}{6}.
        \]
        Check the formula for $n = 1, 2, 3$ before writing the proof.

  \item Prove by strong induction that every natural number $n$ can be written
        in the form $n = 2^{k}m$ with $k \geq 0$ an integer and $m$ an odd
        number. Show where the strong induction hypothesis is used.

  \item Judge the following claims: true, false, or not enough evidence.
        \begin{enumerate}
          \item If $P(1)$ is true and $P(k) \Rightarrow P(k+1)$ for every
                $k \in \mathbb{N}$, then $P(n)$ is true for every $n \in
                \mathbb{N}$.
          \item If $P(n)$ is true for $n = 1, 2, \dots, 10^{6}$, then $P(n)$ is true
                for every $n \in \mathbb{N}$.
          \item If $n^{2}+n+41$ is prime for $n = 0, 1, \dots, 39$, then the claim
                is true for every $n \geq 0$.
          \item If $3$ divides $n$, then $6$ divides $n$.
        \end{enumerate}
        For every item that is false, give one concrete counterexample.

  \item Using Python, count how many primes occur among $n^{2}+n+41$ for
        $n = 0$ up to $n = 100$, then determine all $n$ in that range whose
        value is not prime. Check whether each point of failure has the form
        $n = p-1$ or $n = p$ for some prime $p$.

  \item For $n = 1$ up to $n = 12$, compute $2^{n}$ and $n^{2}$, then determine
        the smallest natural number $n_{0}$ such that $2^{n} > n^{2}$ holds for every
        $n \geq n_{0}$. Prove that claim by induction, with base $n_{0}$.

  \item Prove that the statement ``there is an integer $n > 1$ with
        $n^{2}-1$ prime'' is false. First state the factorization of $n^{2}-1$
        for $n > 1$, then explain why the factors cannot equal $1$.

  \item Show that the statement ``if $n^{2}$ is divisible by $4$ then $n$ is divisible
        by $4$'' is false, then repair the conclusion into a true claim and
        prove it.
\end{enumerate}

\par\smallskip\noindent\textit{Hints for the three hardest items.} Item 5: separate the
case of odd $n$, which is settled at once with $k=0$ and $m=n$, from the case of even $n$,
which uses the induction hypothesis on $n/2$. Item 7: also compute $\sqrt{n^{2}+n+41}$
for every point of failure, and compare it with the smallest divisor. Item 9: $n^{2}-1 = (n-1)(n+1)$; for $n > 1$ both factors lie between $1$
and $n^{2}-1$.
\end{latihan}

\begin{untukpeneliti}
What is genuinely proved in this chapter is three classical facts that survive without a computer:
the irrationality of $\sqrt{2}$ by a parity proof, the infinitude of primes through
Euclid's construction, and the uniqueness of the form of every natural number as a product of primes.
What is merely popularized is the story that $41$ is ``almost'' a prime formula; what happens
is only a long initial stretch, and its length can be predicted from $n^{2}+n+p$.

Primary sources: Abbott, \emph{Understanding Analysis}, Chapter~1 Section~1.1 (irrationality
of $\sqrt{2}$) and Section~1.2 (introduction, including induction); Avigad and Massot,
\emph{Mathematics in Lean}, Chapter~3 (logic) and Chapter~5 Sections~5.1 and 5.2
(irrational roots, induction and recursion); Euclid, \emph{Elements}, Book~IX
Proposition~20 for infinitely many primes.

Open questions: there is no polynomial of positive degree with integer coefficients
that is prime for all $n \geq 0$ (a one-line proof: $f(0)$ divides
$f(t\,f(0))$), but how long a prime stretch a polynomial of high degree can reach
is still being studied, and the distribution of primes is tied to the Riemann hypothesis,
which remains unproved.

Limits of applicability: all results in this chapter are elementary and do not depend on the choice of set
axioms. Formal verification in Lean (Chapter~28) checks the validity of every deductive step,
but it does not check whether the formal statement proved is the one its author
intended.
\end{untukpeneliti}

\peta{Chapter 1 of \emph{Understanding Analysis} (Abbott), in particular Section 1.1 on the
irrationality of $\sqrt{2}$ and Section 1.2 on the introduction; Chapter 3 (logic), Chapter 5
Sections 5.1 and 5.2 (irrational roots as well as induction and recursion) of
\emph{Mathematics in Lean} (Avigad and Massot); and Book IX Proposition 20 of
\emph{Elements} (Euclid) for infinitely many primes.}
